\documentclass[11pt]{article}

\usepackage[margin=1in]{geometry}
\usepackage{booktabs}
\usepackage{pgfplotstable}
\usepackage{graphicx}
\usepackage{hyperref}

\title{ERA5 MIL-HDBK-310 Atmospheric Analysis Report}
\author{Automated MATLAB Pipeline}
\date{\today}

\begin{document}

\maketitle

\section{Introduction}

This report summarizes atmospheric temperature statistics derived from ERA5 reanalysis data using MIL-HDBK-310 exceedance methodology.  
All results are generated automatically from a MATLAB processing pipeline.

\section{Methodology}

Hourly ERA5 temperature data ($t2m$) was processed for multiple cities and time periods.  
Statistics were computed using percentile exceedance levels (1\%, 5\%, 10\%, 90\%, 95\%, 99\%).

Temporal aggregations include:
\begin{itemize}
    \item 1-hour raw resolution
    \item 6-hour moving average
    \item 24-hour moving average
    \item 72-hour moving average
\end{itemize}

\section{Results Summary}

% ==========================================================
% THIS SECTION IS AUTO-GENERATED BY MATLAB
% ==========================================================

%% MATLAB WILL INSERT REPEATED CITY BLOCKS HERE

\section{City-Level Results}
\subsection{AndoyaSpaceCenter}
\subsubsection{Annual MIL-310 Statistics}
\pgfplotstabletypeset[
col sep=comma,
string type
]{AndoyaSpaceCenter/MIL310_summary_AndoyaSpaceCenter_year.csv}


% ---- CITY BLOCK END ----

\section{Conclusion}

The MIL-HDBK-310 statistical framework successfully captures temperature exceedance behavior across multiple temporal scales.  
These results are intended for environmental characterization and engineering design applications.

\end{document}